\begin{pastexam}[label={px:ust-115-ee-ds-41}]{臺聯大電機類 115 學年度 第 41 題}
（單選；2 分）

Answer \textbf{Questions 41-43} according to the following C code.
\begin{minted}[linenos, numbersep=1.5mm, xleftmargin=2em, fontsize=\small]{c}
#include <stdio.h>
int counter; // counting comparison
int binary_search(int arr[], int key, int left, int right) {
    int mid;
    if(left>right){
        counter++;
        return left;
    }
    if (key<arr[left]){ //left boundary
        counter++;
        return left;
    }
    if (key>=arr[right]){ //right boundary
        counter++;
        return right+1;
    }
    mid=(left+right)/2;
    if (key<arr[mid].val){
        counter++;
        binary_search(arr,key,left,mid-1);
    }
    else { // (key>=arr[mid].val)
        binary_search(arr,key,mid+1,right);
    }
}
int main() {
    int arr[9]={1,2,3,5,5,5,6,7,9};
    counter=0;
    int j=binary_search(arr,5,0,8);
    printf("%d",j);
    printf("%d",counter);
}
\end{minted}

What is the output of the code on \textbf{Line 30}?

(A)~2 \qquad (B)~3 \qquad (C)~4 \qquad (D)~5 \qquad (E)~6
\end{pastexam}

\begin{pastsoln}
答案是 (E)。

\textbf{原題的問題}：程式有兩處錯。 第 18、22 行的 \texttt{arr[mid].val}：\texttt{arr} 是 \texttt{int}
陣列，沒有 \texttt{.val} 成員，照原樣編譯不過；第 20、23 行遞迴呼叫沒有寫 \texttt{return}，函式走到
結尾沒有回傳值，\texttt{main} 使用這個回傳值是未定義行為。 本書照明顯的意圖解讀：
\texttt{arr[mid].val} 就是 \texttt{arr[mid]}，兩個遞迴呼叫都是 \texttt{return binary\_search(...)}。

追蹤 \texttt{binary\_search(arr, 5, 0, 8)}，\texttt{arr} $= \langle 1, 2, 3, 5, 5, 5, 6, 7, 9\rangle$
（足標 $0$--$8$）：
\begin{center}
\small
\begin{tabular}{llll}
\toprule
呼叫 & 檢查左右邊界 & 結果 & \texttt{counter} \\
\midrule
$(0, 8)$ & $5 < 1$？否；$5 \ge 9$？否 & $\mathit{mid} = 4$，$5 < 5$？否，往右 & 0 \\
$(5, 8)$ & $5 < 5$？否；$5 \ge 9$？否 & $\mathit{mid} = 6$，$5 < 6$？是，往左 & 1 \\
$(5, 5)$ & $5 < 5$？否；$5 \ge 5$？是 & 回傳 $\mathit{right} + 1 = 6$ & 2 \\
\bottomrule
\end{tabular}
\end{center}
所以 \texttt{j} $= 6$、\texttt{counter} $= 2$。

這個函式回傳的是「第一個大於 \texttt{key} 的位置」：$\texttt{key} \ge \texttt{arr[mid]}$ 時往右找，
所以相等的元素都被跳過。 $\texttt{arr}$ 裡的 $5$ 在足標 $3, 4, 5$，第一個大於 $5$ 的是足標 $6$ 的
$6$，和追蹤一致（同樣的概念見資結書虛擬碼~17.2）。 第 30 行印出 $6$。

其他選項：(A) $2$ 是最後一個小於 $5$ 的位置；(B) $3$ 是第一個等於 $5$ 的位置（「第一個大於等於」
的答案）；(C) $4$ 是 \texttt{mid} 第一次的值；(D) $5$ 是最後一個 $5$ 的位置。 都不是這個函式的回傳值。
\end{pastsoln}
